\documentclass[../../main.tex]{subfiles}

\begin{document}

\section{等价关系、选择公理以及 Zorn 引理}\label{section:ch00-等价关系选择Zorn}

给定两个非空集 $A$ 和 $B$, $A$ 和 $B$ 的笛卡儿积, 记为 $A\times B$, 定义为所有有序对 $(a,b)$ 的族, 其中 $a\in A$ 而 $b\in B$, 且我们考虑 $(a,b)=(a',b')$ 当且仅当 $a=a'$ 且 $b=b'$.\footnote{在基于 Zermelo-Frankel 公理的集合论的正式讨论中, 有序对 $(a,b)$ 定义为集合 $\{\{a\},\{a,b\}\}$, 而具有定义域 $A$ 和值域 $B$ 的函数定义为 $A\times B$ 中有序对的非空族, 它具有以下性质: 若有序对 $(a,b)$ 和 $(a,b')$ 属于该函数, 则 $b=b'$.} 对于非空集合 $X$, 我们称 $X\times X$ 的子集 $R$ 为 $X$ 上的一个关系, 且写作 $xRx'$, 若 $(x,x')$ 属于 $R$. 关系 $R$ 称为自反的, 若对所有 $x\in X$ 有 $xRx$; 关系 $R$ 称为对称的, 若 $x'Rx$ 则 $xRx'$; 关系 $R$ 称为传递的, 若 $xRx'$ 且 $x'Rx''$ 则 $xRx''$.

\begin{definition}[等价关系]\label{definition:ch00-等价关系}
集合 $X$ 上的关系 $R$ 称为等价关系, 若它是自反的、对称的和传递的.
\end{definition}

给定集合 $X$ 上的等价关系 $R$, 对每个 $x\in X$, 集合 $R_{x}=\{x'\mid x'\in X,\,xRx'\}$ 称为 $x$（关于 $R$）的等价类. 等价类族记为 $X/R$. 例如, 给定集合 $X$, 对等关系是 $X$ 的所有子集组成的族 $2^{X}$ 上的等价关系. 一个集合关于对等关系的等价类称为该集合的势.

令 $R$ 为集合 $X$ 上的等价关系. 由于 $R$ 是对称的和传递的, $R_{x}=R_{x'}$ 当且仅当 $xRx'$, 因此等价类族是不交的. 由于关系 $R$ 是自反的, $X$ 是等价类的并. 因此 $X/R$ 是 $X$ 的非空子集的不交族, 其并是 $X$. 反过来, 给定 $X$ 的非空子集的不交族 $\mathcal{F}$, 其并是 $X$, 属于 $\mathcal{F}$ 中的同一个集的关系是 $X$ 上使得 $\mathcal{F}=X/R$ 的等价关系 $R$.

给定集合 $X$ 上的等价关系, 常常有必要选取 $X$ 的子集 $C$, 它恰好由每个等价类的一个成员组成. 这样的集合的存在是显而易见的吗? Ernst Zermelo 唤起了人们对从集族中选取元素这一问题的注意. 比方说, 我们定义两个实数为有理等价, 若它们的差是一个有理数. 容易检验这是实数集上的一个等价关系, 但不易确认一个实数集恰好由每个有理等价类的一个成员组成.

\begin{definition}[选择函数]\label{definition:ch00-选择函数}
令 $\mathcal{F}$ 为非空集的非空簇. $\mathcal{F}$ 上的一个选择函数 $f$ 是从 $\mathcal{F}$ 到 $\bigcup_{F\in\mathcal{F}}F$ 的函数, 它具有以下性质: 对 $\mathcal{F}$ 中的每个集合 $F$, $f(F)$ 是 $F$ 的一个成员.
\end{definition}

\begin{axiom}[Zermelo 选择公理]\label{axiom:ch00-选择公理}
令 $\mathcal{F}$ 为非空集的非空簇, 则 $\mathcal{F}$ 上存在选择函数.
\end{axiom}

非常粗略地说, 非空簇上的选择函数从该簇的每个集合“选取”一个成员. 我们已采用非正式的、描述性的方法引入集合论, 相应地我们将自由地、毫不费力地应用选择公理.

\begin{definition}[偏序、全序、上界与最大元]\label{definition:ch00-偏序}
非空集合 $X$ 上的关系 $R$ 称为偏序, 若它是自反的、传递的, 且对 $X$ 中的 $x,x'$
\begin{align*}
\text{若 }xRx'\text{ 且 }x'Rx,\quad\text{则 }x=x'.
\end{align*}
$X$ 的子集 $E$ 称为是全序的, 若对 $E$ 中的 $x,x'$, 或者 $xRx'$ 或者 $x'Rx$. $X$ 的成员 $x$ 称为是 $X$ 的子集 $E$ 的一个上界, 若对所有 $x'\in E$, $x'Rx$; 而称之为最大的, 若 $X$ 中使得 $x'Rx$ 的唯一成员是 $x'=x$.
\end{definition}

对于集簇 $\mathcal{F}$ 和 $A,B\in\mathcal{F}$, 定义 $ARB$, 若 $A\subseteq B$. 集合的被包含关系是 $\mathcal{F}$ 的偏序. 观察到 $\mathcal{F}$ 中的集合 $F$ 是 $\mathcal{F}$ 的子簇 $\mathcal{F}'$ 的一个上界, 若 $\mathcal{F}'$ 中的每个集合是 $F$ 的子集; 而 $\mathcal{F}$ 中的集合 $F$ 是最大的, 若它不是 $\mathcal{F}$ 中任何集合的真子集. 类似地, 给定集簇 $\mathcal{F}$ 和 $A,B\in\mathcal{F}$, 定义 $ARB$, 若 $B\subseteq A$. 集合的包含关系是 $\mathcal{F}$ 的偏序. 观察到 $\mathcal{F}$ 中的集合 $F$ 是 $\mathcal{F}$ 的子簇 $\mathcal{F}'$ 的一个上界, 若 $\mathcal{F}'$ 中的每个集合包含 $F$; 而 $\mathcal{F}$ 中的集合 $F$ 是最大的, 若它不真包含 $\mathcal{F}$ 中的任何集合.

\begin{lemma}[Zorn 引理]\label{lemma:ch00-zorn}
令 $X$ 为偏序集. 它的每个全序子集有一个上界, 则 $X$ 有一个最大元.
\end{lemma}

我们将用 Zorn 引理证明一些重要的结果, 包括 Hahn-Banach 定理、Tychonoff 乘积定理、Krein-Milman 定理. Zorn 引理等价于 Zermelo 选择公理. 该等价性和相关等价关系的证明, 见 Kelley\cite{Kel75}, pp.\ 31--36.

我们已定义了两个集合的笛卡儿积. 对一般的参数化集簇定义笛卡儿积是有用的. 对于由集合 $\Lambda$ 参数化的集簇 $\{E_{\lambda}\}_{\lambda\in\Lambda}$ 的笛卡儿积, 记为 $\prod_{\lambda\in\Lambda}E_{\lambda}$, 定义为从 $\Lambda$ 到 $\bigcup_{\lambda\in\Lambda}E_{\lambda}$ 使得对每个 $\lambda\in\Lambda$, $f(\lambda)$ 属于 $E_{\lambda}$ 的函数 $f$ 的集合. 显然选择公理等价于非空集的非空簇的笛卡儿积是非空的这一断言. 注意到笛卡儿积是对参数化的集簇定义的, 而相同的簇的两个不同的参数化将有不同的笛卡儿积. 笛卡儿积的这个一般定义与对两个集合给出的定义一致. 事实上, 考虑两个非空集 $A$ 和 $B$. 定义 $\Lambda=\{\lambda_{1},\lambda_{2}\}$, 其中 $\lambda_{1}\neq\lambda_{2}$, 接着定义 $E_{\lambda_{1}}=A$ 与 $E_{\lambda_{2}}=B$. 该映射将有序对 $(f(\lambda_{1}),f(\lambda_{2}))$ 指派给函数 $f\in\prod_{\lambda\in\Lambda}E_{\lambda}$ 是一个将笛卡儿积 $\prod_{\lambda\in\Lambda}E_{\lambda}$ 映到有序对族 $A\times B$ 的可逆映射, 因此这两个集合是对等的. 对于两个集合 $E$ 和 $\Lambda$, 对所有 $\lambda\in\Lambda$ 定义 $E_{\lambda}=E$, 则笛卡儿积 $\prod_{\lambda\in\Lambda}E_{\lambda}$ 等于由所有从 $\Lambda$ 到 $E$ 的映射组成的集合且记为 $E^{\Lambda}$.

\end{document}